\documentclass[answers,12pt,addpoints]{exam} 
\usepackage{import}
\usepackage{hyperref}

\import{C:/Users/prana/OneDrive/Desktop/MathNotes}{style.tex}

% Header
\newcommand{\name}{Pranav Tikkawar}
\newcommand{\course}{Stochastic Processes}
\newcommand{\assignment}{Homework 1}
\author{\name}
\title{\course \ - \assignment}

\begin{document}
\maketitle

\section*{Problem 1: \texorpdfstring{$\sigma$}{sigma}-algebras}

Let $\Omega = \{1,2,3,4,5,6\}$ with the uniform probability measure (a fair die). Define
\[
Z(\omega) =
\begin{cases}
0, & \omega \in \{1,2,3\},\\
1, & \omega \in \{4,5,6\}.
\end{cases}
\]
Let $\mathcal{G} = \sigma(Z)$.

\subsection*{1(a) List all sets in \texorpdfstring{$\mathcal{G}$}{G} and identify its atoms}

$\mathcal{G} = \{ \emptyset, \{1,2,3\}, \{4,5,6\}, \Omega \}$

The atoms of $\mathcal{G}$ are $\{1,2,3\}$ and $\{4,5,6\}$.

\subsection*{1(b) Measurability of \texorpdfstring{$X(\omega) = \mathbf{1}\{\omega \text{ is even}\}$}{X}}

To test if $X$ is $\mathcal{G}$-measurable, we check if the preimage of any set generated by the values of $X$ is in $\mathcal{G}$. The function $X$ takes the value 1 for even $\omega$ (2, 4, 6) and 0 for odd $\omega$ (1, 3, 5).

The preimage of $\{1\}$ under $X$ is $\{2,4,6\}$, which is not in $\mathcal{G}$. Therefore, $X$ is not $\mathcal{G}$-measurable.

\subsection*{1(c) Compute \texorpdfstring{$\mathbb{E}[X \mid \mathcal{G}]$}{E[X|G]}}

To compute $\mathbb{E}[X \mid \mathcal{G}]$, we need to find the conditional expectation of $X$ given the $\sigma$-algebra $\mathcal{G}$. Since $\mathcal{G}$ has two atoms, we compute the expectation on each atom.

\begin{align*}
    E[X \mid \mathcal{G}](\omega) &= \begin{cases}
        E[X \mid Z=0], & \omega \in \{1,2,3\},\\
        E[X \mid Z=1], & \omega \in \{4,5,6\}.
    \end{cases}
    E[X \mid Z=0] &= \frac{1}{3}(X(1) + X(2) + X(3)) = \frac{1}{3}(0 + 1 + 0) = \frac{1}{3},\\
    E[X \mid Z=1] &= \frac{1}{3}(X(4) + X(5) + X(6)) = \frac{1}{3}(1 + 0 + 1) = \frac{2}{3}.
\end{align*}
Thus,
$$\mathbb{E}[X \mid \mathcal{G}](\omega) =
\begin{cases}
\frac{1}{3}, & \omega \in \{1,2,3\},\\
\frac{2}{3}, & \omega \in \{4,5,6\}.
\end{cases}$$


\subsection*{1(d) Compute \texorpdfstring{$\mathbb{E}[Y \mid \mathcal{G}]$}{E[Y|G]} for $Y(\omega) = \omega$}

To compute $\mathbb{E}[Y \mid \mathcal{G}]$, we again compute the expectation on each atom of $\mathcal{G}$.

\begin{align*}
    E[Y \mid Z=0] &= \frac{1}{3}(Y(1) + Y(2) + Y(3)) = \frac{1}{3}(1 + 2 + 3) = 2,\\
    E[Y \mid Z=1] &= \frac{1}{3}(Y(4) + Y(5) + Y(6)) = \frac{1}{3}(4 + 5 + 6) = 5.
\end{align*}

Thus,
$$\mathbb{E}[Y \mid \mathcal{G}](\omega) =
\begin{cases}
2, & \omega \in \{1,2,3\},\\
5, & \omega \in \{4,5,6\}.
\end{cases}$$

\section*{Problem 2: A martingale from prediction (Doob martingale)}

Let $\xi_1,\dots,\xi_6$ be i.i.d.\ Bernoulli$(1/2)$ (1 = head, 0 = tail). Let
\[
\mathcal{F}_n = \sigma(\xi_1,\dots,\xi_n), \qquad
S_n = \sum_{i=1}^n \xi_i.
\]
Define
\[
X = \mathbf{1}\Big\{\sum_{i=1}^6 \xi_i \ge 4\Big\}, \qquad
M_n = \mathbb{E}[X \mid \mathcal{F}_n], \quad n = 0,1,\dots,6.
\]

\subsection*{2(a) Closed-form expression for \texorpdfstring{$M_n$}{Mn} in terms of \texorpdfstring{$S_n$}{Sn}}

This is the conditional expectaion of $X$ given the first $n$ coin flips. Given $\mathcal{F}_n$ we know $S_n$, the number of heads in the first $n$ flips. 

Define
$$ K := \sum_{i=n+1}^6 \xi_i, \sim \text{Binomial}(6-n, 1/2)$$
Thus $\sum_{i=1}^6 \xi_i = S_n + K$.

So then we have
\begin{align*}
    M_n &= \mathbb{P}(S_n + K \geq 4 | \mathcal{F}_n) = \mathbb{P}(K \geq 4 - S_n) \\
    &= \sum_{k = \max(0, 4 - S_n)}^{6-n} \binom{6-n}{k} \left(\frac{1}{2}\right)^{6-n}.
\end{align*}

\subsection*{2(b) Compute \texorpdfstring{$M_0$}{M0}}

We want $M_0 = \mathbb{E}[X \mid \mathcal{F}_0] = \mathbb{E}[X]$. Since $X$ is the indicator of the event that at least 4 heads occur in 6 flips, we can compute this probability directly:
\begin{align*}
M_0 &= \mathbb{P}(S_6 \geq 4) = \sum_{k=4}^6 \binom{6}{k} \left(\frac{1}{2}\right)^6 \\
&= \binom{6}{4} \left(\frac{1}{2}\right)^6 + \binom{6}{5} \left(\frac{1}{2}\right)^6 + \binom{6}{6} \left(\frac{1}{2}\right)^6 \\
&= \left(\frac{1}{2}\right)^6 \left( \binom{6}{4} + \binom{6}{5} + \binom{6}{6} \right) \\
&= \left(\frac{1}{2}\right)^6 (15 + 6 + 1) = \frac{22}{64} = \frac{11}{32}.
\end{align*}

\subsection*{2(c) Compute \texorpdfstring{$M_3$}{M3} on the event $\{\xi_1 = 1, \xi_2 = 0, \xi_3 = 1\}$}

On the event $\{\xi_1 = 1, \xi_2 = 0, \xi_3 = 1\}$, we have $S_3 = 2$. Thus, $S_3 = 2$ and we need to compute
\begin{align*}
M_3 &= \mathbb{P}(S_3 + K \geq 4 | \mathcal{F}_3) = \mathbb{P}(K \geq 4 - S_3) = \mathbb{P}(K \geq 2) \\
&= \sum_{k=2}^{3} \binom{3}{k} \left(\frac{1}{2}\right)^3 = \binom{3}{2} \left(\frac{1}{2}\right)^3 + \binom{3}{3} \left(\frac{1}{2}\right)^3 \\
&= \left(\frac{1}{2}\right)^3 (3 + 1) = \frac{4}{8} = \frac{1}{2}.
\end{align*}

\subsection*{2(d) Prove \texorpdfstring{$(M_n)$}{Mn} is a martingale w.r.t.\ \texorpdfstring{$(\mathcal{F}_n)$}{Fn}}

To show that $(M_n)$ is a martingale with respect to $(\mathcal{F}_n)$, we need to verify the following conditions:
\begin{itemize}
    \item $\mathbb{E}[|M_n|] < \infty$ for all $n$.
    \item $M_n$ is $\mathcal{F}_n$-measurable for all $n$.
    \item $\mathbb{E}[M_{n+1} \mid \mathcal{F}_n] = M_n$ for all $n$.
\end{itemize}
1. Since $M_n$ is a conditional expectation of an indicator random variable, it takes values in $[0,1]$. Thus, $\mathbb{E}[|M_n|] \leq 1 < \infty$.

2. By definition, $M_n = \mathbb{E}[X \mid \mathcal{F}_n]$ is $\mathcal{F}_n$-measurable.

3. We need to show that $\mathbb{E}[M_{n+1} \mid \mathcal{F}_n] = M_n$. Using the tower property of conditional expectation, we have
\begin{align*}
\mathbb{E}[M_{n+1} \mid \mathcal{F}_n] &= \mathbb{E}[\mathbb{E}[X \mid \mathcal{F}_{n+1}] \mid \mathcal{F}_n] \\
&= \mathbb{E}[X \mid \mathcal{F}_n] = M_n.
\end{align*}
Thus, $(M_n)$ is a martingale with respect to $(\mathcal{F}_n)$.

\section*{Problem 3: Wald's lemma and the negative binomial mean}

Let $\xi_1, \xi_2, \dots$ be i.i.d.\ Bernoulli$(p)$ with $p \in (0,1)$. Let
\[
S_n = \sum_{i=1}^n \xi_i, \qquad \mathcal{F}_n = \sigma(\xi_1,\dots,\xi_n).
\]
Fix an integer $r \ge 1$ and define the stopping time
\[
T = \inf\{n \ge 0 : S_n = r\}.
\]

\subsection*{3(a) Show that \texorpdfstring{$T$}{T} is a stopping time w.r.t.\ \texorpdfstring{$(\mathcal{F}_n)$}{Fn}}

A process $T$ is a stopping time with respect to the filtration $(\mathcal{F}_n)$ if for every $n \geq 0$, the event $\{T = n\}$ is in $\mathcal{F}_n$.

Fix $n \geq 0$. The event $\{T = n\}$ can be expressed as
$$ \{T = n\} = \{S_{n} = r\} \cap \{S_k < r \text{ for all } k < n\}. $$

Each $S_k$ is a function of $\xi_1, \dots, \xi_k$, and thus is $\mathcal{F}_k$-measurable. Since $\mathcal{F}_k \subseteq \mathcal{F}_n$ for $k < n$, the event $\{S_k < r\}$ is also in $\mathcal{F}_n$. Therefore, the intersection of these events is in $\mathcal{F}_n$, which means $\{T = n\} \in \mathcal{F}_n$. Since this holds for all $n$, $T$ is a stopping time with respect to $(\mathcal{F}_n)$.

\subsection*{3(b) Show that \texorpdfstring{$Z_n = S_n - pn$}{Zn} is a martingale w.r.t.\ \texorpdfstring{$(\mathcal{F}_n)$}{Fn}}

To show that $Z_n = S_n - pn$ is a martingale with respect to $(\mathcal{F}_n)$, we need to verify the following conditions:
\begin{itemize}
    \item $\mathbb{E}[|Z_n|] < \infty$ for all $n$.
    \item $Z_n$ is $\mathcal{F}_n$-measurable for all $n$.
    \item $\mathbb{E}[Z_{n+1} \mid \mathcal{F}_n] = Z_n$ for all $n$.
\end{itemize}
1. Since $S_n$ is a sum of $n$ Bernoulli random variables, it takes values in $\{0, 1, \dots, n\}$. Thus, $|Z_n| = |S_n - pn| \leq n + pn < \infty$, so $\mathbb{E}[|Z_n|] < \infty$.
2. $Z_n$ is a function of $\xi_1, \dots, \xi_n$, so it is $\mathcal{F}_n$-measurable.
3. We need to show that $\mathbb{E}[Z_{n+1} \mid \mathcal{F}_n] = Z_n$. We have
\begin{align*}
\mathbb{E}[Z_{n+1} \mid \mathcal{F}_n] &= \mathbb{E}[S_{n+1} - p(n+1) \mid \mathcal{F}_n] \\
&= \mathbb{E}[S_n + \xi_{n+1} - p(n+1) \mid \mathcal{F}_n] \\
&= S_n - pn + \mathbb{E}[\xi_{n+1} \mid \mathcal{F}_n] - p \\
&= S_n - pn + p - p = Z_n.
\end{align*}
Thus, $Z_n$ is a martingale with respect to $(\mathcal{F}_n)$.

\subsection*{3(c) Apply bounded Wald's lemma to \texorpdfstring{$Z_{T_N}$}{Z\_{T\_N}}}

Define $T_N = T \wedge N$.
Since $T_N$ is a bounded stopping time (it is at most $N$), we can apply the bounded version of Wald's lemma to the martingale $(Z_n)$ at the stopping time $T_N$.

By Wald's lemma, we have
$$
\mathbb{E}[Z_{T_N}] = \mathbb{E}[Z_0] 
$$
Since $Z_0 = S_0 - p \cdot 0 = 0$ and $Z_{T_N} = S_{T_N} - p T_N$, we get
\begin{align*}
    \mathbb{E}[Z_{T_N}] &= \mathbb{E}[S_{T_N} - p T_N] = E[S_{T_N}] - p \mathbb{E}[T_N] = 0 \\
    \Rightarrow \mathbb{E}[S_{T_N}] &= p \mathbb{E}[T_N].
\end{align*}
Thus we have shown that $\mathbb{E}[S_{T_N}] = p \mathbb{E}[T_N]$.

\subsection*{3(d) Take the limit \texorpdfstring{$N \to \infty$}{N to infty} and conclude \texorpdfstring{$\mathbb{E}[T] = r/p$}{E[T]=r/p}}

As $N \to \infty$, the stopping time $T_N = T \wedge N$ converges to $T$ almost surely. Since $S_{T_N} = r$ for all $N$ large enough (specifically, for all $N \geq T$), we have
$$
\mathbb{E}[S_{T_N}] = r.
$$
Taking the limit as $N \to \infty$ in the equation $\mathbb{E}[S_{T_N}] = p \mathbb{E}[T_N]$, we get
$$
r = p \lim_{N \to \infty} \mathbb{E}[T_N] = p \mathbb{E}[T].
$$
Thus, we conclude that
$$
\mathbb{E}[T] = \frac{r}{p}.
$$



\end{document}